%% kapa: exam=1 round=36 year=2025 period=1 form=A subject=real_estate subject_ko=부동산학원론 seq=11 number=91
%% source: Q-Net 감정평가사 기출문제 (cst00309) artlSeq=5231525 file=raw/1008/5231525/2025년 제36회 감정평가사 1차시험 1교시 문제지 원본.pdf
%% conversion: kordoc 4.24.51 -> tools/gen1.py (auto); review=proofread-sonnet; flags=none
\begin{question}{91}
여과과정과 주거분리에 관한 설명으로 옳지 않은 것은? (단, 주어진 조건에 한함)
\begin{choices}
  \item 여과과정이란 시간의 흐름에 따라 특정 주택의 질적 변화와 외부성이 복합적으로 작용해 주택가치가 변하게 되면서 상이한 소득계층들의 전ㆍ출입이 진행되는 것을 말한다.
  \item 고소득층 주거지역에 인접한 저소득층 주거지역에서 주택개량을 통한 가치상승분이 주택개량비용보다 작은 경우, 상향여과과정이 발생한다.
  \item 상향여과과정은 소득증가 등의 이유로 인해 저가주택의 수요가 감소할 때 나타날 수 있다.
  \item 주거분리현상은 지리적으로 인접한 근린지역에서 뿐만 아니라 도시 전체에서도 발생할 수 있다.
  \item 침입과 계승의 현상으로 인해 주거입지의 변화가 나타날 수 있다.
\end{choices}
\end{question}
\answer{2}
